\documentclass[../../../include/open-logic-section]{subfiles}

\begin{document}

\olfileid[fr]{sth}{replacement}{refproofs}
\olsection{Annexe : résultats autour du remplacement}

Dans cette section, nous démontrerons le schéma de réflexion
dans $\ZF$. Nous établirons aussi en quel sens la réflexion
équivaut au remplacement. Enfin, nous démontrerons une conséquence
intéressante de ces résultats concernant la force de la
réflexion et du remplacement. \emph{Attention : il s'agit
de loin des mathématiques les plus avancées de ce manuel.}

Commençons par un lemme qui, pour alléger la notation, emploie
une \emph{barre suscrite} pour désigner des suites de variables
ou d'objets. Ainsi, «~$\overline{a}_k$~» abrège
«~$a_{k_1}$, \dots, $a_{k_n}$~», où $n$ est déterminé
par le contexte.

\begin{lem}\ollabel{lemreflection}
Pour chaque $1 \leq i \leq k$, soit
$\phi_i(\overline{v}_i, x)$ une formule. Alors, pour tout
$\alpha$, il existe un $\beta > \alpha$ tel que,
quels que soient $\overline{a}_1, \ldots,
\overline{a}_k \in V_\beta$ et quel que soit
$1 \leq i \leq k$ :
\[
	\exists x\phi_i(\overline{a}_i, x) \rightarrow
	(\exists x \in V_\beta) \phi_i(\overline{a}_i, x)
\]
\end{lem}

\begin{proof}
Définissons un terme $\mu$ comme suit :
$\mu(\overline{a}_1, \ldots, \overline{a}_k)$
est la première étape $V$ qui satisfait, pour
$1 \leq i \leq k$, toutes les implications suivantes :
\[
\exists x\phi_i(\overline{a}_i, x) \rightarrow
(\exists x \in V) \phi_i(\overline{a}_i, x)
\]
On vérifie aisément que
$\mu(\overline{a}_1, \ldots, \overline{a}_k)$
existe quels que soient $\overline{a}_1, \ldots,
\overline{a}_k$. En utilisant le remplacement et notre
théorème de récursion, définissons :
\begin{align*}
	S_0 & = V_{\alpha+1}\\
	S_{n+1} & = S_n \cup \bigcup
	\Setabs{\mu(\overline{a}_1, \ldots, \overline{a}_k)}
	{\overline{a}_1, \ldots, \overline{a}_k \in S_n} \\
	S &= \bigcup_{m < \omega} S_m.
\end{align*}
Chaque $S_n$, et donc $S$ lui-même, est une étape
postérieure à $V_\alpha$.\footnote{Note de l'édition française :
dans l'original, la condition définissant $\mu$ comporte
une parenthèse fermante en trop, et la dernière ligne de la
définition de $S$ porte $S_n$ au lieu de $S_m$ sous l'indice
$m$. Ces deux coquilles sont corrigées ici.}
Fixons $\overline{a}_1$, \dots,~
$\overline{a}_k \in S$ ; il existe donc un
$n < \omega$ tel que $\overline{a}_1$, \dots,
$\overline{a}_k \in S_n$. Fixons $1 \leq i \leq k$
et supposons $\exists x\phi_i(\overline{a}_i,x)$.
Par construction,
$(\exists x \in \mu(\overline{a}_1,
\ldots, \overline{a}_k))\phi_i(\overline{a}_i, x)$.
Il s'ensuit que
$(\exists x \in S_{n+1})\phi_i(\overline{a}_i, x)$,
et donc que $(\exists x \in S)\phi_i(\overline{a}_i, x)$.
Ainsi, $S$ est le $V_\beta$ recherché.
\end{proof}
\noindent
Nous pouvons maintenant démontrer le
\olref[sth][replacement][ref]{reflectionschema}
assez directement :

\begin{proof}[Démonstration]
Fixons $\alpha$. Sans perte de généralité, nous pouvons
supposer que les seuls opérateurs logiques de $\phi$ sont
$\exists$, $\lnot$ et $\land$ (puisqu'ils suffisent
à tout exprimer). Énumérons les sous-formules de $\phi$
par ordre de complexité : $\psi_1, \ldots, \psi_k$,
avec $\psi_k = \phi$. D'après le \olref{lemreflection},
il existe un $\beta > \alpha$ tel que, quels que soient
$\overline{a}_i \in V_\beta$ et $1 \leq i \leq k$ :
\begin{align}\label{reflectionnicelybehaved}
	\exists x\psi_i(\overline{a}_i, x) \rightarrow
	(\exists x \in V_\beta) \psi_i(\overline{a}_i, x)\tag{*}
\end{align}
Par récurrence sur la complexité de $\psi_i$, montrons que
$\psi_i(\overline{a}_i) \leftrightarrow
\psi_i^{V_\beta}(\overline{a}_i)$ pour tout
$\overline{a}_i \in V_\beta$. Si $\psi_i$ est atomique,
c'est immédiat. Si $\psi_i$ est une négation ou une
conjonction de sous-formules qui vérifient cette
propriété, $\psi_i$ la vérifie également.
Le seul cas intéressant est donc celui de
la quantification. Fixons
$\overline{a}_i \in V_\beta$ ; alors :
\begin{align*}
	(\exists x \psi_i(\overline{a}_i, x))^{V_\beta}
	&\text{ ssi }
	(\exists x \in V_\beta)\psi_i^{V_\beta}(\overline{a}_i, x)
	&&\text{par définition}\\
	&\text{ ssi }
	(\exists x \in V_\beta)\psi_i(\overline{a}_i,  x)
	&&\text{par hypothèse}\\
	&\text{ ssi }
	\exists x \psi_i(\overline{a}_i, x)
	&&\text{par \eqref{reflectionnicelybehaved}}
\end{align*}
La récurrence est ainsi achevée ; le résultat suit de
$\psi_k = \phi$.
\end{proof}

Nous avons démontré la réflexion dans $\ZF$.
Notre démonstration suit pour l'essentiel
\citet{Montague1961}. Nous voulons maintenant
montrer, dans $\Z$, que la réflexion entraîne
le remplacement. La démonstration suit
\citet{Levy1960}, moyennant une simplification.

Puisque nous travaillons dans $\Z$, nous ne pouvons
pas présenter la réflexion exactement sous la forme
ci-dessus. Nous l'avons en effet formulée au moyen de
la notation «~$V_\alpha$~», qui ne peut être définie
dans $\Z$ (voir la \olref[sth][spine][zf]{sec}).
Nous proposerons donc à la place la formulation
apparemment plus faible suivante de la réflexion :
\footnote{Note de l'édition française : l'original dit ici
«~du remplacement~», mais la définition qui suit est celle
de la réflexion faible.}

\begin{defish}
\emph{Réflexion faible.} Pour toute formule $\phi$,
il existe un ensemble transitif $S$ tel que $0$, $1$
et tous les paramètres de $\phi$ soient des !!{element}s
de $S$, et
$(\forall \overline{x} \in S)(\phi \liff \phi^S)$.
\end{defish}

Pour en déduire le remplacement, suivons d'abord
\citet[première partie du théorème 2]{Levy1960} et montrons
que nous pouvons «~réfléchir~» deux formules à la fois :

\begin{lem}[dans $\Z + \text{Réflexion faible}$]\ollabel{lem:reflect}
Pour toutes formules $\psi, \chi$, il existe un ensemble
transitif $S$ tel que $0$ et $1$ (ainsi que tous les
paramètres des formules) soient des !!{element}s de $S$,
et que
$(\forall \overline{x} \in S)((\psi \liff \psi^S)
\land (\chi \liff \chi^S))$.
\end{lem}

\begin{proof}
Soit $\phi$ la formule
$(z = 0 \land \psi) \lor (z = 1 \land \chi)$, où $z$ ne figure pas librement dans $\psi$ ni dans $\chi$.

Nous employons ici une abréviation : il faudrait écrire
«~$z = 0$~» sous la forme
«~$\forall t\, t \notin z$~», et «~$z =1$~»
sous la forme
«~$\forall s(s \in z \liff \forall t\, t \notin s)$~».
Mais puisque $0, 1 \in S$ et que $S$ est transitif,
ces formules sont \emph{absolues} pour $S$ :
elles s'appliquent au même objet, que l'on restreigne
ou non leurs quantificateurs à $S$.\footnote{Plus
formellement, pour $z \in S$, si $\xi$ est l'une de ces deux formules,
$\xi(z) \liff \xi^S(z)$.}

Par réflexion faible, nous obtenons un ensemble $S$
approprié tel que :
\begin{align*}
	(\forall z, \overline{x} \in S)(&\phi \liff \phi^S)\\
	\text{soit }(\forall z, \overline{x} \in S)(&((z = 0 \land \psi) \lor (z = 1 \land \chi)) \liff {}\\
	&\phantom{(}((z = 0 \land \psi) \lor (z = 1 \land \chi))^S)\\
	\text{soit }(\forall z, \overline{x} \in S)(&((z = 0 \land \psi) \lor (z = 1 \land \chi))\liff {}\\
	&\phantom{(}((z = 0 \land \psi^S) \lor (z = 1 \land \chi^S)))\\
	\text{soit }(\forall \overline{x} \in S)(&(\psi \liff \psi^S) \land (\chi \liff \chi^S))
\end{align*}
La deuxième assertion entraîne la troisième parce que
«~$z = 0$~» et «~$z=1$~» sont absolues pour $S$ ;
la quatrième suit de $0 \neq 1$.
\end{proof}\noindent
Nous pouvons maintenant obtenir le remplacement en
suivant et en simplifiant
\citet[théorème 6]{Levy1960} :

\begin{thm}[dans $\Z$ + Réflexion faible]\label{thm:replacement}
Pour toute formule $\phi(v,w)$ et tout $A$, si
$(\forall x \in A)\lexists![y][\phi(x,y)]$, alors
l'ensemble $\Setabs{y}{(\exists x \in A)\phi(x,y)}$
existe.\footnote{Note de l'édition française : il manque
une parenthèse fermante dans l'expression de l'ensemble
dans l'original ; elle est rétablie ici.}
\end{thm}

\begin{proof}
Fixons $A$ tel que
$(\forall x \in A)\lexists![y][\phi(x,y)]$,
et définissons les formules :
\begin{align*}
	\psi &\text{ est } (\phi(x, z) \land A = A)\\
	\chi &\text{ est } \lexists[y][\phi(x, y)]
\end{align*}
Puisque $A$ est un paramètre de $\psi$, le
\olref{lem:reflect} fournit un ensemble
transitif~$S$ tel que $0, 1, A \in S$ (ainsi
que tout autre paramètre), et tel que :
\[
	(\forall x,z \in S)((\psi \liff \psi^S)
	\land (\chi \liff \chi^S))
\]
En particulier :
\begin{align*}
	(\forall  x, z \in S)(&\phi(x, z) \liff \phi^S(x, z))\\
	(\forall x \in S)(&\exists y\phi(x, y) \liff
	(\exists y \in S)\phi^S(x, y))
\end{align*}
En combinant ces énoncés et en remarquant que
$A \subseteq S$, puisque $A \in S$ et que
$S$ est transitif, on obtient :
\begin{align*}
	(\forall x \in A)(&\exists y\phi(x, y) \liff
	(\exists y \in S)\phi(x, y))
\end{align*}
Or $(\forall x \in A)(\lexists![y \in S])\phi(x, y)$,
car $(\forall x \in A)\lexists![y][\phi(x, y)]$.
La séparation fournit donc
$\Setabs{y \in S}{(\exists x \in A) \phi(x, y)}
= \Setabs{y}{(\exists x \in A) \phi(x, y)}$.
\end{proof}

\end{document}
